% 与questions_ch文件配套；来源超链接见sources.tex。
\Answer{SUPL5-01}{组间差异的变化}
选C。预设对比为
\[
(\bar Y_{2,4}-\bar Y_{1,4})-(\bar Y_{2,2}-\bar Y_{1,2})=(10-8)-(3-4)=3.
\]
它估计这一特定时间对比上的组间差异变化。仅由这些均数不能给出其标准误或检验结果。
\par\nobreak\smallskip\nobreak{\footnotesize\color{Muted}\noindent 来源与改编：案例改编。材料及改编范围见\hyperref[src:L5-UCLA-CONTRAST]{[22]}、\hyperref[src:L5-COURSE]{[21]}。\par}

\Answer{SUPL5-02}{随机截距与个体内差值}
选B。差值中的$b_i$相消，故方差为$4+4=8$。亦可由$\operatorname{Var}(Y_{ij})=13$、$\operatorname{Cov}(Y_{i4},Y_{i1})=9$算得$13+13-2\times9=8$。该模型的协方差为CS结构，是讲义HF结构的特例。
\par\nobreak\smallskip\nobreak{\footnotesize\color{Muted}\noindent 来源与改编：知识点配套。题干与设问为新编；参考依据\hyperref[src:L5-COURSE]{[21]}。\par}

\Answer{SUPL5-03}{与基线比较的多重性}
选A。排序后的阈值依次为$0.05/3$、$0.05/2$、$0.05$。第1项通过，第2项$0.028>0.025$，顺序检验在此停止。Holm校正后的$P$值为$0.036,\ 0.056,\ 0.056$。第3项不能越过前面的未拒绝项单独通过。
\par\nobreak\smallskip\nobreak{\footnotesize\color{Muted}\noindent 来源与改编：知识点配套。题干与设问为新编；参考依据\hyperref[src:L5-COURSE]{[21]}、\hyperref[src:L5-R-HOLM]{[23]}。\par}

\Answer{SUPL5-04}{不等间隔的变化率对比}
选D。所需指标为
\[
\frac{Y_2-Y_1}{2-1}-\frac{Y_{0.5}-Y_0}{0.5-0}
=2Y_0-2Y_{0.5}-Y_1+Y_2.
\]
分母应采用实际经过的小时数；用访视序号代替实际时间会改变“变化率”的含义。
\par\nobreak\smallskip\nobreak{\footnotesize\color{Muted}\noindent 来源与改编：案例改编。材料及改编范围见\hyperref[src:L5-RABBIT]{[24]}、\hyperref[src:L5-COURSE]{[21]}。\par}

\Answer{SUPL5-05}{SEM能说明什么}
选E。$S=0.64\sqrt{20}\approx2.86$，因原表已四舍五入，换算值也是近似值。3次测量仍来自同一批婴儿；同一人的两次测量通常相关，变化量的标准误还需要协方差。因此边际SEM不能单独确定变化量或交互作用的检验。
\par\nobreak\smallskip\nobreak{\footnotesize\color{Muted}\noindent 来源与改编：案例改编。材料及改编范围见\hyperref[src:L5-RUEDA]{[27]}、\hyperref[src:L5-COURSE]{[21]}。\par}

\Answer{SUPL5-06}{正交对比与独立性}
选D。I中的$Z_i$为相互独立的单个正态标量，直接估计其方差即可。II错误：两对比的协方差为$c_1^\top\Sigma c_2$，一般不能用$c_1^\top c_2$替代。III由$C^\top\Sigma C=\lambda I$及多元正态分布下不相关蕴含独立得到。这里的HF结构是广义球形性，不要求原始$\Sigma=\sigma^2I$。
\par\nobreak\smallskip\nobreak{\footnotesize\color{Muted}\noindent 来源与改编：知识点配套。题干与设问为新编；参考依据\hyperref[src:L5-COURSE]{[21]}。\par}

\Answer{SUPL5-07}{漏访与实际时间不一致}
删除不完整对象会损失信息，并可能因退出与病情有关而造成选择偏倚；不能把完整病例自动视为原随机分组的代表。按统一访视名称作分析还可能掩盖实际经过时间的差异，尤其当研究目标是单位时间的变化率时。

应保留对象编号、组别、计划与实际测量时间、已观测结局及退出原因；先描述缺失模式，再考虑使用线性混合效应模型处理不平衡测量，并按研究目标把实际时间作为连续变量或合理建模。该方法仍依赖均数和协方差模型；常用似然推断还依赖给定已观测信息后可忽略缺失的条件（如MAR）。若退出还依赖未观测结局，应作相应敏感性分析。不能笼统声称“混合模型自动消除缺失偏倚”。计划时点本身不等间隔，并不使把时间作为分类因素的RM ANOVA自动无效。
\par\nobreak\smallskip\nobreak{\footnotesize\color{Muted}\noindent 来源与改编：知识点配套。题干与设问为新编；参考依据\hyperref[src:L5-COURSE]{[21]}、\hyperref[src:L5-UCLA-SAS]{[26]}。\par}

\Answer{SUPL5-08}{重复测量输出的误差分层}
(1) 个体（组内）误差均方为$24.25/6=4.041667$，个体内误差均方为$76/12=6.333333$。因此
\[
F_{\text{组}}=\frac{2542.041667}{4.041667}=628.959,\quad
F_{\text{时}}=\frac{1/2}{6.333333}=0.07895,\quad
F_{\text{交互}}=\frac{1736.333333/2}{6.333333}=137.079.
\]
参考分布分别为$F_{1,6}$、$F_{2,12}$、$F_{2,12}$。组效应用个体之间的误差；时间及组别$\times$时间用个体内误差，不能共用错误的分母。

(2) 不能。总体时间效应反映跨组平均的时间变化，组内变化可能互相抵消。显著交互提示两组时间模式不同，不能把时间主效应不显著解释成两组都不变。平方和和$F$值不带方向，仅凭此表不能判断哪组上升，也不能把显著组主效应解释为每个时点均有组间差异。

(3) 可根据预设科学问题比较特定时点的组间简单效应，或比较两组从基线到某时点的变化量，也可设定其他有意义的组$\times$时间对比；应报告估计值及区间。对预先界定的一组后续比较，采用Holm等方法控制多重性，避免只挑显著的比较。
\par\nobreak\smallskip\nobreak{\footnotesize\color{Muted}\noindent 来源与改编：案例改编。材料及改编范围见\hyperref[src:L5-UCLA-SAS]{[26]}、\hyperref[src:L5-COURSE]{[21]}、\hyperref[src:L5-R-HOLM]{[23]}。\par}

\Answer{SUPL5-09}{比较预设的0至2小时变化量}
(1) 第1组变化量为$17.2,-4.4,-48.0,-33.0,-35.0$，第3组为$228.1,48.3,62.4,70.5,87.3$。故
\[
\bar D_1=-20.64,\quad s_1^2=700.788;\qquad
\bar D_3=99.32,\quad s_3^2=5381.372.
\]
(2) 差异估计为$99.32-(-20.64)=119.96$，
\[
SE=\sqrt{\frac{5381.372}{5}+\frac{700.788}{5}}=34.8774,\qquad
t=\frac{119.96}{34.8774}=3.4395,\quad \nu\approx5.0244.
\]
双侧$P\approx0.0183$。

(3) 使用未舍入的临界值，95\%置信区间为$(30.44,\ 209.48)$；使用题给$2.567$得近似$(30.43,\ 209.49)$，均应接受。在给定模型假定下，第3组的平均0至2小时变化比第1组大，区间较宽，估计仍不精确。

这一检验利用每只兔的实际差值保留了配对信息，检验的是一个预设的组间变化量对比。它不等于3组4时点的总体交互作用检验，也不能单独说明0.5或1小时的组间差异。
\par\nobreak\smallskip\nobreak{\footnotesize\color{Muted}\noindent 来源与改编：案例改编。材料及改编范围见\hyperref[src:L5-RABBIT]{[24]}、\hyperref[src:L5-RABBIT-TRACE]{[25]}、\hyperref[src:L5-COURSE]{[21]}。\par}

\Answer{SUPL5-10}{重复测量文献的分析目标与结局定义}
(1) 没有。逐时点组间检验和组内时间检验分别针对不同假设，并不直接检验两组的变化模式是否相同。“一组显著、另一组不显著”也不等于变化差异显著。应在符合结局分布与个体相关结构的方法中，明确检验组$\times$时间项或预设的组间变化量对比。

(2) 不足。边际SEM不含同一婴儿各次测量的协方差。至少应取得可按婴儿编号连接的个体纵向记录、组别、实际时点和缺失情况；若随机化确有匹配或分层，还应保留相应设计信息。取得数据后再检查适用假定并选择方法，不能从表II拼出交互作用的$P$值。原文因分布检查结果而采用非参数方法，不能未经审查直接替换为正态RM ANOVA。

(3) 不能。跨时点比较同时改变了阈值，所定义的二分类结局并不完全相同。宜采用可比的固定阈值，或直接分析保留尺度信息的计数/对数计数及其时间变化。同一时点两组使用相同阈值，因此不存在这个跨时点阈值不一致的问题；但仍需按研究设计和比较范围解释。

(4) 完成病例分析可能因退出与结局或组别相关而偏倚；仅凭3人退出也不能确定偏倚方向。应按随机组报告人数流转、原因和已观测结局，说明分析人群；可结合退出机制对缺失结局进行适当敏感性分析。不能把完成者结果自动称为全部随机入组对象的分析。
\par\nobreak\smallskip\nobreak{\footnotesize\color{Muted}\noindent 来源与改编：案例改编。材料及改编范围见\hyperref[src:L5-RUEDA]{[27]}、\hyperref[src:L5-COURSE]{[21]}。\par}
